\begin{tikzpicture}[font=\small,>=stealth,line join=round]
\draw[draw=teal!45!black,fill=green!3,rounded corners=3pt] (0,0) rectangle (13.6,6.5);
\node at (6.8,5.95) {定义域即使左右对称，中间断开时仍不能直接反推};
\begin{scope}[shift={(4.7,1.7)},x=1.2cm,y=0.58cm]
\draw[gray,->] (-2.7,0)--(2.7,0) node[right] {$x$};
\draw[gray,->] (0,-0.3)--(0,6.5);
\fill[gray!12] (-1,-.25) rectangle (1,6);
\draw[teal!75!black,very thick,domain=-2.2:-1,samples=100] plot(\x,{(\x)^2+1});
\draw[teal!75!black,very thick,domain=1:2.2,samples=100] plot(\x,{(\x)^2});
\draw[teal!75!black,fill=green!3,thick] (-1,2) circle(2.5pt) (1,1) circle(2.5pt);
\fill[teal!75!black] (-2,5) circle(2.5pt);
\fill[teal!75!black] (2,4) circle(2.5pt);
\node[above] at (-2,5.2) {$5$};\node[above] at (2,4.2) {$4$};
\draw[gray] (-2,0+.08)--(-2,0-.08);\node[below] at (-2,0-.15) {$-2$};
\draw[gray] (-1,0+.08)--(-1,0-.08);\node[below] at (-1,0-.15) {$-1$};
\draw[gray] (1,0+.08)--(1,0-.08);\node[below] at (1,0-.15) {$1$};
\draw[gray] (2,0+.08)--(2,0-.08);\node[below] at (2,0-.15) {$2$};
\node at (0,6.6) {缺少中间区间};
\end{scope}
\node[anchor=west,align=left,text width=4.2cm] at (8.7,3.2) {$D=(-\infty,-1)\cup(1,\infty)$。\\[8pt]左边 $f=x^2+1$，\\右边 $f=x^2$。\\[8pt]两边都满足 $f^{\prime}=2x$，\\但 $f(-2)=5\ne4=f(2)$。};
\node at (6.8,.4) {两段可以相差不同常数；反推时需要在同一个对称区间上可导。};
\end{tikzpicture}
